\documentclass[11pt,a4paper]{article}
\usepackage[T1]{fontenc}
\usepackage{lmodern,amsmath,amssymb,amsthm,microtype,booktabs}
\usepackage[margin=25mm]{geometry}
\usepackage[colorlinks=true,urlcolor=blue,linkcolor=blue]{hyperref}
\newtheorem{proposition}{Proposition}
\newtheorem{theorem}{Theorem}
\newcommand{\dd}{\mathop{}\!\mathrm d}
\title{The worldline action and uniform local\\Barnes asymptotics in the recoil limit}
\author{Alan Day \and Chengheng\thanks{AI research assistant (OpenAI).}}
\date{Working research note --- 30 October 2026}
\begin{document}\maketitle
\begin{abstract}
We supply an analytic action identity for the reflected recoil
Barnes saddle at general positive Euclidean tetrahedron masses.
An elementary quadratic-field identity identifies each logarithmic
mass derivative with the corresponding exterior dihedral angle;
homogeneity then gives the worldline action. We also establish
uniform estimates for the actual localized Gamma quotient,
retaining the transverse energies. These estimates justify the
integrated removal of its exact periodic pole factor even when the
pole-row separation tends to zero. The result concerns a specified
local contribution. Dominance in the full Barnes cycle, transverse
tails, and the physical four-point dictionary remain open.
\end{abstract}

\section{Definitions and the order of limits}
The de Sitter radius is one. We first take the fixed-mass
\(\kappa\to\infty\) limit, and only afterwards increase the masses.
This does not assert a joint finite-\(G\) limit. The object remains
the internal chiral amplitude of the fixed-topology CLP graph
\cite{CLP}; no new physical state or external energy cycle is specified.

For edge labels \(\mu_1,\ldots,\mu_6\), let \(a\) be the four
half-sums on \(123,156,246,345\), and \(r\) the three half-sums on
\(1245,1346,2356\). The exact reflected kernel is
\begin{align}
 \mathcal B(z)&=H(z)E(z),\\
 H(z)&=8\pi^3e^{2\pi iz-\pi\sum r}
 \frac{\prod_a\Gamma(z+ia)}
 {\Gamma(z+\tfrac12)\prod_r\Gamma(z+\tfrac12+ir)},&
 E(z)&=\prod_r(1+e^{2\pi i(z+ir)})^{-1}.
 \label{eq:exact}
\end{align}
Use \(f(t)=t(\log t-1)\) and
\begin{align}
 \Phi(z)&=2\pi iz-\pi\sum r+
 \sum_a f(z+ia)-\sum_r f(z+ir)-f(z),\\
 g(z)&=8\pi^3\prod_a(z+ia)^{-1/2}.
\end{align}
The branches below are fixed by the right saddle. With \(s=iz\),
\[
 P(s)=\prod_a(s-a)-s\prod_r(s-r)=A_qs^2+Bs+C,\qquad
 \Delta=B^2-4A_qC=-9V_{\rm dual}^2.
\]
Here \(A_q\) is a quadratic coefficient, distinct from the large
parameter \(K\) used later.

\section{An action identity for all six real masses}
\begin{proposition}
For a positive, nondegenerate Euclidean dual tetrahedron, set
\(T=\sqrt{-\Delta}=3V_{\rm dual}>0\),
\(s_+=(-B+iT)/(2A_q)\), and \(z_+=-is_+\).
Then, with the indicated logarithm branches,
\begin{equation}
 -2\Re\Phi(z_+)=\sum_{j=1}^6\mu_jd_j=I_{\rm CSTY},
 \label{eq:action}
\end{equation}
where \(d_j\in(0,\pi)\) are the exterior dihedral angles in
CSTY Appendix B.1 \cite{CSTY}.
\end{proposition}

\begin{proof}
Write \(\boldsymbol\mu=(l,b,c,f,d,e)\).
The four faces are \((l,b,c),(l,d,e),(b,f,e),(c,f,d)\).
The edge opposite \(l\) is \(f\), and
\(A_q=(lf+bd+ce)/2>0\), so \(\Re z_+>0\).
Define
\begin{align}
 D_l&=[l^2-(b-c)^2][l^2-(d-e)^2]>0,\\
 H_l&=(l^2+b^2-c^2)(l^2+e^2-d^2)
       -2l^2(b^2+e^2-f^2).
\end{align}
Let
\[
 a_1=\frac{l+b+c}{2},\quad a_2=\frac{l+d+e}{2},\quad
 r_1=\frac{l+b+f+d}{2},\quad r_2=\frac{l+c+f+e}{2}.
\]
The quadratic relation gives
\begin{equation}
 R_l:=\frac{(z_++ia_1)(z_++ia_2)}
              {(z_++ir_1)(z_++ir_2)}
       =\frac{H_l+8ilT}{D_l}.
 \label{eq:angle-ratio}
\end{equation}
This can be checked without radicals: replace \(iT\) by \(2A_qs+B\)
and cross-multiply; the resulting polynomial is divisible by \(P(s)\).
The new code verifies the zero remainder with all six masses symbolic.
The further polynomial identity
\[
 H_l^2+64l^2T^2=
 [16\,\mathcal A(l,b,c)\mathcal A(l,d,e)]^2
\]
uses the Heron area \(\mathcal A\).
Since \(D_l,l,T>0\), \(R_l\) lies in the upper half-plane.
Its argument has exactly the cosine specified by CSTY (B.6),
hence \(\arg R_l=d_l\).

All arguments \(z_++ia,z_++ir,z_+\) lie in the right half-plane.
The exponentiated saddle equation therefore makes \(\Phi'(z_+)\)
a continuous integer multiple of \(2\pi i\). The domain of
nondegenerate Euclidean edge lengths is connected: squared edge
lengths are linear functions of the positive definite Gram matrix
of the three edges at a chosen vertex.
At the regular tetrahedron the integer is zero, so \(\Phi'(z_+)=0\)
throughout this domain.

Differentiating on shell now gives
\begin{equation}
 \partial_l[-2\Re\Phi(z_+)]
 =2\pi+\Im\left[
 \log(z_++ia_1)+\log(z_++ia_2)
 -\log(z_++ir_1)-\log(z_++ir_2)\right]=d_l .
\end{equation}
The last branch integer is again fixed at the regular tetrahedron
and cannot jump on the connected domain.
The other five edges follow by symmetry.
Finally, \(\sum a=\sum r\) cancels every scaling term
\(\lambda\log\lambda\); the on-shell potential is homogeneous of
degree one. Euler's identity gives \eqref{eq:action}.
\end{proof}

At complex masses, continue the two chiral functions separately:
\[
 \Phi^\dagger(\boldsymbol\mu)
 =\overline{\Phi(\overline{\boldsymbol\mu})},\qquad
 I=-\Phi-\Phi^\dagger .
\]
The equality extends in any common analytic continuation domain
with the transported angle and saddle branches. It does not mean
taking a new modulus at arbitrary complex masses, and it does not
determine which saddle contributes to an integration cycle.

\section{Uniform estimates for the actual local amplitude}
Use the equal-probe labels
\[
 (\mu_1,\ldots,\mu_6)=(m_1,\nu,m_2,m_3,\nu,m_4),\quad
 m_{1,3}=M+\omega/2\pm u,\quad m_{2,4}=M-\omega/2\pm v.
\]
We take the joint large-mass regime
\begin{equation}
 K=\frac{\nu^2}{2M}\to\infty,\quad \epsilon=\frac{\nu}{M}\to0,\quad
 \omega=iKy,\quad |y|\le Y,\quad
 \rho=\frac{|u|+|v|}{\nu}\to0,\qquad u,v\in\mathbb R.
 \label{eq:regime}
\end{equation}
This includes transverse boxes
\(|u|,|v|\le L\sqrt M\), \(L=o(\sqrt K)\).
One can impose the earlier WKB window
\(\nu^{4/3}\ll M\ll\nu^2\) in addition.
The local result below does not extend that separate WKB theorem.

Write \(z_+=x_+-ic\). Direct quadratic algebra gives
\begin{equation}
 \delta:=M+\nu-c=
 \frac{\nu(\nu^2-u^2-v^2+\omega^2)}
 {4M^2+2\nu^2-2u^2-2v^2+\omega^2}
 \sim\frac{\nu^3}{4M^2}>0,\qquad
 \frac{x_+}{K}=\frac{\sqrt{1+y^2}}2+O(\epsilon^2+\rho^2).
 \label{eq:gap}
\end{equation}
Choose a fixed
\[
 0<\eta<\frac{\sqrt{1+Y^2}-Y}{4}.
\]
For real \(|h|\le\eta K\), the two near denominator parameters are
\[
 z_++h+i(M+\nu\pm\omega/2)
 =x_++h\mp Ky/2+i\delta,
\]
whose real parts exceed \(c_YK\).
The four numerator arguments have modulus comparable to \(\nu\);
the two other denominator arguments have modulus comparable to \(M\).
All stay in a fixed sector away from the negative real axis.
The left pole rows remain above the line and the right rows below.
Thus the horizontal contour can follow this imaginary-energy path
without pole crossing. This finite-parameter statement is not a
global asymptotic estimate.

Let \(\chi\) be a fixed even \(C^4\) cutoff, \(0\le\chi\le1\),
equal to one for \(|t|\le1/2\) and zero for \(|t|\ge1\), and define
\begin{equation}
 F_K(x)=\chi\left(\frac{x-x_+}{\eta K}\right)
 \frac{H(x-ic)}{g(z_+)e^{\Phi(z_+)}} .
 \label{eq:F}
\end{equation}

\begin{theorem}[Localized Gamma quotient]
In \eqref{eq:regime}, uniformly in the stated parameter sets,
\begin{align}
 \|F_K^{(4)}\|_1&\le CK^{-3/2},\label{eq:derivative}\\
 \int_{\mathbb R}F_K(x)\dd x
 &=\sqrt{\frac{2\pi K}{q_*}}\,[1+O(K^{-1})],
 &q_*&=-K\Phi''(z_+)
   =4\sqrt{1+y^2}+O(\epsilon+\rho^2).
 \label{eq:gaussian}
\end{align}
The square root is chosen with \(\Re q_*>0\).
In particular the integral has modulus at least \(c\sqrt K\).
\end{theorem}

\begin{proof}
The sector and modulus bounds permit the fixed-shift Stirling
expansion with remainder \cite{DLMF}:
\[
 H(z)=g(z)e^{\Phi(z)+R(z)},\qquad
 R^{(j)}(z)=O(K^{-1-j})\quad(0\le j\le4).
\]
For each local point, a complex disk of radius \(c'K\) stays in
a safe sector for all Gamma arguments. The scalar \(O(K^{-1})\)
remainder and Cauchy's estimates give the derivative bounds.
The half-integer shifts in \eqref{eq:exact} are already included
in \(g\), and
\[
 (\log g)^{(j)}=O(\nu^{-j})\quad(1\le j\le4).
\]
Direct differentiation also gives
\(\Phi^{(j)}=O(K^{1-j})\), \(2\le j\le6\).
The two near denominators contribute
\(-1/(z+ir_+)-1/(z+ir_-)\) to \(\Phi''\).
Their real part is at most \(-c/K\);
all other terms are \(O(\nu^{-1})=o(K^{-1})\).
Consequently throughout the local segment
\[
 \Re[\Phi(z_++h)-\Phi(z_+)]\le -ch^2/(2K).
\]
Here stationarity is exact, by continuation of \(\Phi'=0\).

Apply the fourth exponential derivative formula to
\(\Phi+\log g+R\), then differentiate the cutoff.
With constants adjusted as necessary, the result is bounded by
\[
 CK^{-2}(1+|h|/\sqrt K)^4e^{-c'h^2/K}
 +CK^{-4}e^{-c''K}
   \mathbf1_{\{\eta K/2\le|h|\le\eta K\}} .
\]
Integration proves \eqref{eq:derivative}.
The cutoff terms have not been dropped; their support is
exponentially suppressed.

For the integral, put \(h=\sqrt K\,t\).
The quadratic exponent is \(-q_*t^2/2\).
The cubic phase and the linear amplitude correction are odd
polynomials of size \(O(K^{-1/2})\), so their integrals against
the even complex Gaussian and even cutoff vanish.
Fourth-order terms, the remainder \(R\), and squares of the odd
corrections give \(O(K^{-1})\).
For a uniform justification first expand on \(|t|\le K^{1/12}\);
the preceding concavity estimate controls its exterior.
Finally \(q_*\) lies in a fixed compact subset of the right
half-plane and has the displayed limit, as follows from
\(\Phi''=-\sqrt{-\Delta}/\prod_a(z_++ia)\).
The Gaussian integral is bounded away from zero, proving the
relative estimate and lower bound.
\end{proof}

\section{The pole factor is now controlled inside the local integral}
The preceding note proved the following elementary estimate:
for \(|q_r|<1\), \(E_q(x)=\prod_{r=1}^3(1+q_re^{2\pi ix})^{-1}\)
and \(F\in W^{4,1}(\mathbb R)\),
\begin{equation}
 \left|\int F(E_q-1)\right|\le
 C_4\|F^{(4)}\|_1,\qquad
 C_4=\frac{\zeta(2)+3\zeta(3)+2\zeta(4)}{2(2\pi)^4}.
 \label{eq:fourier}
\end{equation}
Indeed the \(n\)-th Fourier coefficient is bounded by
\(\binom{n+2}{2}\); four integrations by parts give a summable bound.
The constant is independent of how close the \(q_r\)'s are to the
unit circle, although the endpoint itself is excluded.

On our line \(q_r=e^{-2\pi(r-c)}\) and \(|q_r|<1\).
The theorem verifies both hypotheses needed to turn
\eqref{eq:fourier} into the relative estimate
\begin{equation}
 \frac{\left|\int F_K(E-1)\right|}
 {\left|\int F_K\right|}=O(K^{-2}),
 \label{eq:poleerror}
\end{equation}
uniformly even when \(\delta\downarrow0\).
Thus the actual local contribution
\[
 \mathfrak B_{\rm loc}=\frac1{2\pi i}\int_{\mathbb R}
 \chi\left(\frac{x-x_+}{\eta K}\right)
 \mathcal B(x-ic)\dd x
\]
obeys
\begin{equation}
 \mathfrak B_{\rm loc}=
 \frac{g(z_+)e^{\Phi(z_+)}}
 {i\sqrt{2\pi[-\Phi''(z_+)]}}\,[1+O(K^{-1})].
 \label{eq:localresult}
\end{equation}
The orientation is left to right; the Gaussian branch is fixed by
\(\Re[-\Phi'']>0\).
This establishes the local expression using the actual Gamma
quotient and pole factor. It does not replace
\(\mathfrak B_{\rm loc}\) by the full \(\mathfrak B\).
Real \(\omega\) away from zero is also outside the uniform regime
of this theorem: its proposed saddle line can cross pole rows.

\section{Diagnostics and the remaining global problem}
The script \texttt{code/recoil\_local.py} records 50 passing checks:
general symbolic angle, modulus and Cayley--Menger identities, the
positive quadratic coefficient and transverse pole height, ten cutoff jets,
all six angles on three unequal tetrahedra, two separate-chiral
complex continuations, and refined local quadratures.
The proofs above, rather than finite-point sampling, support
the uniform claims.

For the local calculation take \(\epsilon=K^{-3/2}\),
\(u=.3\sqrt M\), \(v=-.2\sqrt M\), and \(\eta=.1\).
The cutoff transition is the degree-nine polynomial
\(1-126s^5+420s^6-540s^7+315s^8-70s^9\), with
\(s=2|h|/(\eta K)-1\) on its transition interval.
It has the required four vanishing endpoint jets.
The actual Gamma quotient, rather than its Gaussian approximation,
is integrated.
\begin{center}
\begin{tabular}{rrrr}
\toprule
\(K\)&\(y\)&pole height \(\delta\)&relative Gaussian error\\
\midrule
400&0&\(2.4996\times10^{-2}\)&\(3.7970\times10^{-3}\)\\
1600&1&\(1.2499\times10^{-2}\)&\(9.1941\times10^{-6}\)\\
6400&1&\(6.2499\times10^{-3}\)&\(2.3011\times10^{-6}\)\\
\bottomrule
\end{tabular}
\end{center}
Quadrature is refined from 48 to 96 nodes per segment at 45/60 digits.
The complex integral is separately repeated at 80 digits on the
same 96-node segmentation.
The fourth-derivative \(L^1\) norm has nearly sharp absolute-value
corners. The initial unsplit grid failed its fixed tolerance;
splitting at zeros of the actual derivative's real part resolved it.
Both the initial code and failed output are retained.
These are empirical diagnostics, not certified interval errors.
We have not directly quadrature-tested the complete large-mass
integral with \(E\); its local control follows from
\eqref{eq:fourier}--\eqref{eq:poleerror}.

The next required result is a relative bound on
\(\mathfrak B-\mathfrak B_{\rm loc}\), including the other saddle,
the negative-real direction and both remote tails.
Transverse exterior regions, physical insertion factors, the three
external energy cycles and disconnected baseline remain separate
requirements. The work does not establish finite-\(G\) dynamics,
a microscopic state dictionary, or a complete de Sitter dual.
This is an author self-reviewed proof draft without independent
certification or a priority claim.

\begin{thebibliography}{9}
\bibitem{CSTY} Y. Chen, D. Stanford, H. Tang and Z. Yang,
\emph{Negative shocks versus static patch holography},
\href{https://arxiv.org/html/2607.14042v1}{arXiv:2607.14042v1},
Appendix B.1.
\bibitem{CLP} J. Chandra, D. Liska and B. Post,
\emph{3D de Sitter from complex Virasoro TQFT: Observers and knots},
\href{https://arxiv.org/html/2610.30247v1}{arXiv:2610.30247v1}.
\bibitem{DLMF} NIST DLMF,
\href{https://dlmf.nist.gov/5.11}{Section 5.11},
fixed-sector Stirling expansions, shifted forms and remainder bounds.
\end{thebibliography}
\end{document}
